Put \(q=m-2c=h-c\). By \(\text{def:affine-separation}\), \(q<3\) would imply \(\gamma_H=0\), contradicting \(\gamma_H\geq\gamma_0>0\). Hence
\[
q\geq3.
\tag{1}
\]

\textbf{Uniform matrix bounds.}
Let \(G\) denote either the true structural matrix or a feasible candidate matrix satisfying the theorem's cycle-slack and condition-number bounds. In the Leibniz expansion of \(\det G\), the identity-permutation term equals one because \(G\) has unit diagonal. Every other permutation decomposes into fixed points and disjoint directed simple cycles. Every nontrivial cycle product in \(I_p-G\) has absolute value at most \(1-\zeta_0\in[0,1)\), while every fixed-point factor has absolute value one. Thus each of the \(p!\) Leibniz terms has absolute value at most one, and
\[
|\det G|\leq p!.
\tag{2}
\]
Write \(\sigma_1(G)=\|G\|_{\mathrm{op}}\) and order the singular values decreasingly. The condition-number bound implies \(\sigma_j(G)\geq\sigma_1(G)/\overline\kappa\) for \(2\leq j\leq p\). Therefore
\[
\frac{\sigma_1(G)^p}{\overline\kappa^{p-1}}
\leq\prod_{j=1}^p\sigma_j(G)
=|\det G|
\leq p!,
\]
and hence
\[
\|G\|_{\mathrm{op}}\leq L_0.
\tag{3}
\]
Since a unit diagonal entry gives \(\|G\|_{\mathrm{op}}\geq1\),
\[
\|G^{-1}\|_{\mathrm{op}}
=\frac{\|G\|_{\mathrm{op}}\|G^{-1}\|_{\mathrm{op}}}
{\|G\|_{\mathrm{op}}}
\leq\overline\kappa.
\tag{4}
\]

\textbf{A uniform compact residual gap.}
Set
\[
\rho_0=\min\left\{1,\frac{1}{2K_0\overline M}\right\}>0.
\tag{5}
\]
For each choice of sets \(H_1,H_2\subseteq E\) of size \(h\) and \(S\subseteq H_1\cap H_2\) of size \(q\), consider tuples
\[
\tau=\bigl(D,U,\Omega,(s_e,\Sigma_e)_{e\in H_1},
A,V,\Psi,(t_e,\Gamma_e)_{e\in H_2}\bigr)
\]
subject to the following closed conditions: \(DU=UD=I_p\) and \(AV=VA=I_p\); both structural matrices have unit diagonal, cycle slack at least \(\zeta_0\), and condition number at most \(\overline\kappa\); \(\Omega,\Psi,\Sigma_e,\Gamma_e\succeq0\), and \(s_e,t_e\in[0,\infty)^p\); each side satisfies its BACKSHIFT covariance equations; each side's affine separation, computed over every \(q\)-element subset of its size-\(h\) fitted set, is at least \(\gamma_0\); and each side satisfies the displayed scale bound with upper bound \(\overline M\). Let \(\mathfrak T\) be the finite union of these tuple families.

The inverse equations and (3)--(4) bound \(D,U,A,V\). Because every term in either scale inequality is nonnegative, the scale bounds separately bound the positive-semidefinite nuisances, all covariance matrices, and every nonnegative shift coordinate. Every remaining constraint is closed: the cycle conditions and affine-minor conditions are finite weak polynomial inequalities, the positive-semidefinite cone and nonnegative orthant are closed, the norm bounds are closed by continuity, and the covariance and inverse relations are polynomial equalities. Thus every fixed-index piece is closed and bounded in a finite-dimensional Euclidean space. The elementary finite-dimensional Heine--Borel theorem makes each piece compact, and therefore the finite union \(\mathfrak T\) is compact. This compactness argument uses positive semidefiniteness and the scale bounds, but no covariance inverse or covariance eigenvalue lower bound.

Define the continuous residual
\[
\mathfrak r(\tau)=
\max_{e\in S}
\left\|A\Sigma_eA^{\mathsf T}
-\Psi-\operatorname{diag}(t_e)\right\|_{\mathrm F}
\]
and the closed far family
\[
\mathfrak T_{\mathrm{far}}
=\{\tau\in\mathfrak T:\|AU-I_p\|_{\mathrm F}\geq\rho_0\}.
\]
If \(\mathfrak T_{\mathrm{far}}=\varnothing\), set \(\mu_0=1\). Otherwise compactness and continuity give an attained minimum
\[
\mu_0=\min_{\tau\in\mathfrak T_{\mathrm{far}}}\mathfrak r(\tau).
\tag{6}
\]
We show that this minimum is positive. If a far tuple had \(\mathfrak r(\tau)=0\), then on \(S\)
\[
D\Sigma_eD^{\mathsf T}=\Omega+\operatorname{diag}(s_e),
\qquad
A\Sigma_eA^{\mathsf T}=\Psi+\operatorname{diag}(t_e).
\tag{7}
\]
For each \(k<\ell\), the first side's affine-separation condition supplies three labels in \(S\) whose affine minor has absolute value at least \(\gamma_0\). Hence the first side's \((k,\ell)\)-shift cloud on \(S\) is not contained in an affine line.

Because \(|E\setminus S|=m-q=2c\), partition \(E\setminus S\) into disjoint sets \(C_1,C_2\) of cardinality \(c\), and define a common covariance family by
\[
\widetilde\Sigma_e=
\begin{cases}
\Sigma_e,&e\in S,\\
D^{-1}\Omega D^{-\mathsf T},&e\in C_1,\\
A^{-1}\Psi A^{-\mathsf T},&e\in C_2.
\end{cases}
\]
The first explanation fits \(S\cup C_1\) using its shifts on \(S\) and zero shifts on \(C_1\); the second fits \(S\cup C_2\) using its shifts on \(S\) and zero shifts on \(C_2\). The added covariances are positive semidefinite because \(\Omega,\Psi\succeq0\) and congruence by an invertible matrix preserves positive semidefiniteness. The added shifts are nonnegative. Both fitted sets have cardinality \(q+c=h\), their intersection is \(S\), and (7) makes their covariance equations agree on that intersection. All hypotheses of \(\text{lem:overlap-uniqueness}\) therefore hold, so \(A=D\). Since \(U=D^{-1}\), this gives \(AU=I_p\), contradicting membership in \(\mathfrak T_{\mathrm{far}}\). Thus
\[
\mu_0>0.
\tag{8}
\]
This also holds under the convention \(\mu_0=1\) when the far family is empty. Define, non-effectively,
\[
r_0=\frac{\mu_0}{2\sqrt p\,L_0^2}>0.
\tag{9}
\]
Because \(\mathfrak T\), and hence \(\mu_0\), was defined solely from the fixed dimensions and the four numerical bounds, \(r_0\) has exactly the asserted uniform dependence. No decision or computation of \(\mu_0\) is used.

\textbf{Entry into the identity neighborhood.}
Fix an outcome in \(G_{n,\alpha,H}\), let \(r=r_{n,\alpha,H}\leq r_0\), and take any \(A\in\mathcal U^{(n)}_\alpha\). By \(\text{def:confidence-union}\), select \(\Gamma_e\in\mathcal C^{(n)}_{e,\alpha}\) for all \(e\in E\) such that \(A\in\mathcal R_c(\Gamma)\). By \(\text{def:compatible-set}\), there are a fitted set of size at least \(h\), a nuisance \(\Psi\succeq0\), and nonnegative shifts satisfying the candidate covariance equations. Restrict the fitted set and its witnesses to a subset \(K\subseteq E\) of size exactly \(h\). The theorem's declarative candidate hypothesis supplies all four candidate-side bounds on this restriction.

Inclusion--exclusion gives
\[
|H\cap K|\geq|H|+|K|-|E|=2h-m=q.
\]
Choose \(S\subseteq H\cap K\) with \(|S|=q\). Put \(U=D^{-1}\) and \(V=A^{-1}\). The true and candidate equations and their respective four bounds show that the resulting objects form a tuple in \(\mathfrak T\). Since the outcome is covered and \(e\in S\subseteq H\),
\[
\|\Gamma_e-\Sigma_e\|_{\mathrm{op}}\leq r
\qquad(e\in S).
\tag{10}
\]
The candidate covariance equation, (3), (10), submultiplicativity, and \(\|X\|_{\mathrm F}\leq\sqrt p\|X\|_{\mathrm{op}}\) imply
\[
\begin{aligned}
\left\|A\Sigma_eA^{\mathsf T}
-\Psi-\operatorname{diag}(t_e)\right\|_{\mathrm F}
&=\left\|A(\Sigma_e-\Gamma_e)A^{\mathsf T}\right\|_{\mathrm F}\\
&\leq\sqrt p\,L_0^2r
\leq\mu_0/2.
\end{aligned}
\tag{11}
\]
If the tuple belonged to \(\mathfrak T_{\mathrm{far}}\), then in the nonempty-far case (6) would make the maximum in (11) at least \(\mu_0\), a contradiction; in the empty-far case membership is impossible by definition. Consequently, with
\[
Q=AD^{-1}=I_p+R,
\qquad u=\|R\|_{\mathrm F},
\]
one has
\[
u<\rho_0\leq1,
\qquad K_0\overline M u\leq\frac12.
\tag{12}
\]

\textbf{Local affine-minor inversion.}
Fix \(e_0\in S\) and write \(x_e=s_e-s_{e_0}\). For every \(i<j\), the true-side affine-separation bound supplies a triple \(a,b,d\in S\) whose \((i,j)\) affine determinant has absolute value at least \(\gamma_0\). Writing \(\widetilde x_e=(x_{ei},x_{ej})\in\mathbb R^2\), translation invariance of the affine determinant and the identity
\[
\det(\widetilde x_a-\widetilde x_b,
\widetilde x_d-\widetilde x_b)
=\det(\widetilde x_a,\widetilde x_d)
+\det(\widetilde x_d,\widetilde x_b)
+\det(\widetilde x_b,\widetilde x_a)
\]
show, by the triangle inequality, that two labels \(e,f\in S\) satisfy
\[
|x_{ei}x_{fj}-x_{ej}x_{fi}|\geq\gamma_0/3.
\tag{13}
\]
Let \(E_e=\Gamma_e-\Sigma_e\). Subtract the candidate equation at \(e_0\) from that at \(e\), and subtract the true covariance equation at \(e_0\) from that at \(e\). The common nuisances cancel. Taking the \((i,j)\) off-diagonal entry yields
\[
R_{ij}x_{ej}+R_{ji}x_{ei}
=-\sum_{k=1}^pR_{ik}R_{jk}x_{ek}
-[A(E_e-E_{e_0})A^{\mathsf T}]_{ij}.
\tag{14}
\]
The scale and nonnegativity conditions give \(|x_{ek}|\leq\overline M\). Equation (10) gives \(\|E_e-E_{e_0}\|_{\mathrm{op}}\leq2r\), and
\[
\|A\|_{\mathrm{op}}=\|QD\|_{\mathrm{op}}
\leq(1+u)L_0.
\]
Consequently, Cauchy--Schwarz and \(\sum_k|R_{ik}R_{jk}|\leq u^2\) turn (14) into
\[
|R_{ij}x_{ej}+R_{ji}x_{ei}|
\leq\overline M u^2+2rL_0^2(1+u)^2.
\tag{15}
\]
Apply (15) at the two labels from (13) and solve the resulting two-by-two linear system for \(R_{ij}\) and \(R_{ji}\). Each coefficient has absolute value at most \(\overline M\), so Cramer's rule and (13) give, for every ordered pair \(i\neq j\),
\[
|R_{ij}|
\leq\frac{6\overline M}{\gamma_0}
\left\{\overline M u^2+2rL_0^2(1+u)^2\right\}.
\tag{16}
\]
The unit diagonals of \(A\) and \(D\) imply
\[
0=(A-D)_{ii}=(RD)_{ii}
=R_{ii}+\sum_{k\neq i}R_{ik}D_{ki}.
\tag{17}
\]
For each \(i\), the squared norm of the \(i\)-th column of \(D\) is at most \(L_0^2\). Squaring (16), summing over the \(p(p-1)\) ordered off-diagonal entries, applying Cauchy--Schwarz to (17), and taking square roots therefore yield
\[
u\leq K_0
\left\{\overline M u^2+2rL_0^2(1+u)^2\right\}.
\tag{18}
\]
By (12), the quadratic term on the right is at most \(u/2\), and \((1+u)^2\leq4\). Hence
\[
u\leq16K_0L_0^2r.
\tag{19}
\]
It follows that
\[
\|A-D\|_{\mathrm{op}}
=\|RD\|_{\mathrm{op}}
\leq uL_0
\leq16K_0L_0^3r
=C_0r.
\tag{20}
\]
Taking the supremum over \(A\in\mathcal U^{(n)}_\alpha\) proves the outer-radius bound. On \(G_{n,\alpha,H}\), \(\text{thm:confidence-union-coverage}\) gives \(D\in\mathcal U^{(n)}_\alpha\), so the union is nonempty. For \(A_1,A_2\in\mathcal U^{(n)}_\alpha\), the triangle inequality and (20) give
\[
\|A_1-A_2\|_{\mathrm{op}}
\leq\|A_1-D\|_{\mathrm{op}}+\|A_2-D\|_{\mathrm{op}}
\leq2C_0r,
\]
which proves the diameter bound. Because the same \(r_0\) and \(C_0\) apply throughout the fixed four-margin class, \(r_{n,\alpha,H}\to0\) on covered outcomes implies both asserted limits.